%============================================================ 2. 관련 연구
\section{관련 연구}
\label{sec:related}

\subsection{무기-표적 할당과 조합 최적화}

방어 자산을 표적에 배분하는 문제는 무기-표적 할당(WTA)으로 정식화되어 정확해법과
휴리스틱이 모두 연구되었고~\citep{ahuja2007wta}, 1:1 배정으로 축소하면 헝가리안
알고리즘~\citep{kuhn1955hungarian}으로 다항 시간에 최소비용 매칭을 얻는다. 이 방법들은
비용 함수가 주어지면 최적을 보장하지만, ``앞서 놓은 그물이 이미 막은 회랑은 두고 남은
그물을 열린 쪽으로 돌려라''처럼 이전 결정의 결과와 자산 잔량에 따라 달라지는 판단은 고정된
거리·각도 비용으로 표현되지 않는다. 본 연구는 배정 계층에 조합 최적화 대신 언어모델의
판단을 두어 전술적 맥락을 반영하고, 그 대가를 \S\ref{sec:results}에서 정량적으로 보고한다.

\subsection{무인 위협 대응}

무인 위협의 무력화 수단은 연성(전자전)·경성(물리 타격)·포획(그물 등)으로
구분된다~\citep{wang2021cuas}. 유선 조종 표적의 확산~\citep{sutea2026fiber}은 연성 대응의
범위를 좁혔고, 요격 비용의 비대칭~\citep{rumbaugh2024cost}은 경성 대응의 지속성을 제약한다.
해상에서는 다중 USV 협력 기동~\citep{liu2016usv}과 다중 에이전트 강화학습에 의한 협조
포위~\citep{zhang2025encircle}가 보고되었으나, 포위는 방어정이 표적보다 빠르거나 대등함을
전제한다. 본 연구는 방어정이 \emph{더 느린} 조건에서 추격 대신 접근 경로의 선제 차단을 택한다.

\subsection{다중 에이전트 강화학습과 크레딧 배분}

중앙 집중 학습·분산 실행(CTDE) 계열은 팀 보상에서 개체별 기여를 분리하기 위해 가치망을
학습한다 --- 중앙 비평가(MADDPG~\citep{lowe2017maddpg}), 단조 가치 분해(QMIX~\citep{rashid2018qmix}),
반사실적 기준선(COMA~\citep{foerster2018coma}), PPO~\citep{schulman2017ppo}의 다중 에이전트
확장(MAPPO~\citep{yu2022mappo}). 공유 가치망의 추정 편향과 학습 불안정~\citep{haarnoja2018sac}은
다른 에이전트가 함께 학습하는 비정상 환경에서 더 커진다~\citep{lowe2017maddpg}. 언어모델 정렬 분야의 그룹 상대 정책 최적화(GRPO)~\citep{shao2024deepseekmath}는
같은 상태에서 뽑은 여러 후보의 보상을 서로 비교해 이득을 산정함으로써 가치망을 없앤다.
본 연구는 GRPO를 다중 에이전트 제어로 옮기고 후보 하나를 기하 휴리스틱으로 고정해 기준선으로
삼으며, 후보의 팀 보상을 비평가가 아니라 \emph{시뮬레이터 재진행}으로 얻는다. COMA의
반사실적 착상을 같은 방식으로 구현한 변형도 비교하였다.

\subsection{격자 지도 표현과 픽셀 지목 행동}

완전 합성곱 구조~\citep{long2015fcn}와 U-Net~\citep{ronneberger2015unet}은 모든 픽셀에 값을
매기는 조밀 예측을 가능하게 하고, CoordConv~\citep{liu2018coordconv}는 좌표 채널로 절대
위치를, FiLM~\citep{perez2018film}은 조건 벡터에 의한 채널별 변조를 더한다. 이들은 모두
상황 \emph{지각}에 쓰였다. 후보마다 점수를 매겨 하나를 \emph{지목}하는 출력은 포인터
네트워크~\citep{vinyals2015pointer}에서 비롯되었고, 그 점수는 어텐션~\citep{vaswani2017attention}과
같은 질의--후보 내적이지만 주로 조합 최적화에 쓰였다. 본 연구는 둘을 잇는다 --- U-Net 특징
지도의 각 픽셀에 방어정별 질의로 점수를 매기고(\S\ref{sec:policy}), 그 점수맵을 \emph{정책의
행동 분포}로 삼아 픽셀 하나를 지목하는 것이 곧 행동이다. 격자 행동공간이므로 안전 제약을 마스크로
행동 분포에 직접 부과할 수 있다.

\subsection{언어모델 기반 계획}

언어모델을 계획 주체로 쓰는 연구는 고수준 지시를 실행 가능한 단계로 분해하는 방향으로
발전했고~\citep{huang2022zeroshot}, SayCan~\citep{ahn2022saycan}은 사전 지식에 로봇의 실행
가능성 추정을 결합해 불가능한 계획을 걸러낸다. 구조적 출력~\citep{openai2026function}은
언어모델 출력을 하위 시스템이 파싱 가능한 형태로 고정한다. 본 연구는 언어모델
판단만으로 배정이 성립하는지, 어느 능력 수준부터 이득인지를 측정한다.
